\documentclass[a4paper,12pt]{article}
\usepackage[utf8]{inputenc}
\usepackage[T1]{fontenc}
\usepackage[french]{babel}
\usepackage{hyperref}
\usepackage{listings}
\usepackage{xcolor}
\usepackage{graphicx}
\usepackage{geometry}

\geometry{hmargin=2.5cm,vmargin=2.5cm}

\hypersetup{
    colorlinks=true,
    linkcolor=blue,
    filecolor=magenta,      
    urlcolor=cyan,
    pdftitle={Guide de Déploiement - Yowyob Feedback Frontend},
    pdfpagemode=FullScreen,
}

\definecolor{codegreen}{rgb}{0,0.6,0}
\definecolor{codegray}{rgb}{0.5,0.5,0.5}
\definecolor{codepurple}{rgb}{0.58,0,0.82}
\definecolor{backcolour}{rgb}{0.95,0.95,0.92}

\lstdefinestyle{mystyle}{
    backgroundcolor=\color{backcolour},   
    commentstyle=\color{codegreen},
    keywordstyle=\color{magenta},
    numberstyle=\tiny\color{codegray},
    stringstyle=\color{codepurple},
    basicstyle=\ttfamily\footnotesize,
    breakatwhitespace=false,         
    breaklines=true,                 
    captionpos=b,                    
    keepspaces=true,                 
    numbers=left,                    
    numbersep=5pt,                  
    showspaces=false,                
    showstringspaces=false,
    showtabs=false,                  
    tabsize=2
}

\lstset{style=mystyle}

\title{\textbf{Guide de Déploiement : Yowyob Feedback Frontend}}
\author{Équipe Yowyob}
\date{\today}

\begin{document}

\maketitle

\section{Introduction}
Ce document décrit la procédure détaillée pour déployer l'application frontend \textbf{Yowyob Feedback} sur la plateforme \href{https://vercel.com}{Vercel}. L'application est développée avec \textbf{Next.js 16} (App Router), TypeScript et Tailwind CSS.

\section{Prérequis}
Avant de commencer, assurez-vous de disposer des éléments suivants :
\begin{itemize}
    \item Un compte \href{https://github.com}{GitHub} avec accès au dépôt du code source.
    \item Un compte \href{https://vercel.com}{Vercel} (peut être lié à votre compte GitHub).
    \item Les clés API et variables d'environnement nécessaires (Supabase, Backend API, Groq).
\end{itemize}

\section{Configuration de l'Environnement}
L'application nécessite plusieurs variables d'environnement pour fonctionner correctement en production. Voici la liste des variables à configurer sur Vercel :

\begin{itemize}
    \item \texttt{NEXT\_PUBLIC\_SUPABASE\_URL} : URL de votre instance Supabase.
    \item \texttt{NEXT\_PUBLIC\_SUPABASE\_ANON\_KEY} : Clé publique anonyme de Supabase.
    \item \texttt{NEXT\_PUBLIC\_API\_URL} : URL de l'API Backend publique (ex: \texttt{https://api.monsite.com/api/v1}). \\ \textit{Note : Ne mettez pas \texttt{localhost} en production.}
    \item \texttt{GROQ\_API\_KEY} : Clé API pour l'assistant YowBot (Groq).
\end{itemize}

\section{Procédure de Déploiement}

\subsection{Étape 1 : Préparation du Code}
Assurez-vous que votre code est à jour et poussé sur la branche principale (\texttt{main} ou \texttt{master}) de votre dépôt GitHub.
\begin{lstlisting}[language=bash]
git add .
git commit -m "Prêt pour le déploiement"
git push origin main
\end{lstlisting}

\subsection{Étape 2 : Importation dans Vercel}
\begin{enumerate}
    \item Connectez-vous à votre tableau de bord Vercel.
    \item Cliquez sur le bouton \textbf{"Add New..."} puis \textbf{"Project"}.
    \item Sélectionnez votre fournisseur Git (GitHub) et recherchez le dépôt \texttt{yowyob\_feedback\_frontend}.
    \item Cliquez sur \textbf{"Import"}.
\end{enumerate}

\subsection{Étape 3 : Configuration du Projet}
Dans la section \textbf{"Configure Project"} :
\begin{enumerate}
    \item \textbf{Framework Preset} : Vercel devrait détecter automatiquement \textbf{Next.js}.
    \item \textbf{Root Directory} : Laissez la valeur par défaut \texttt{./} (racine du dépôt) sauf si votre projet est dans un sous-dossier.
    \item \textbf{Build and Output Settings} : Les commandes par défaut sont généralement correctes :
    \begin{itemize}
        \item Build Command: \texttt{next build}
        \item Install Command: \texttt{npm install} (ou détecté automatiquement)
    \end{itemize}
\end{enumerate}

\subsection{Étape 4 : Ajout des Variables d'Environnement}
Dépliez la section \textbf{"Environment Variables"} et ajoutez les clés listées dans la section \textbf{Configuration de l'Environnement} (Supabase, API URL, Groq).

\subsection{Étape 5 : Lancement du Déploiement}
Cliquez sur le bouton \textbf{"Deploy"}. Vercel va :
\begin{enumerate}
    \item Cloner votre dépôt.
    \item Installer les dépendances.
    \item Lancer le build de l'application.
    \item Déployer le résultat sur leur réseau Edge.
\end{enumerate}

\section{Vérification}
Une fois le déploiement terminé, Vercel vous fournira une URL (ex: \texttt{https://yowyob-feedback.vercel.app}).
\begin{itemize}
    \item Accédez à cette URL pour vérifier que la page d'accueil se charge correctement.
    \item Testez l'authentification et l'appel à l'assistant YowBot pour confirmer que les variables d'environnement sont bien prises en compte.
\end{itemize}

\end{document}
